\documentclass[10pt,a4paper]{amsart}
\usepackage[margin=19mm]{geometry}
\usepackage{amsmath,amssymb,mathtools}
\usepackage[scr=rsfs,cal=euler]{mathalfa}
\usepackage{xfrac}
\usepackage[unicode,hidelinks,bookmarks=false]{hyperref}
\newcommand{\ie}{i.e.\ }
\newcommand{\cf}{cf.\ }
\newcommand{\resp}{respectively\ }
\newcommand{\Subsection}{\subsection}
\providecommand{\nobreakdash}{\nobreak\hbox{-}}
\setlength{\emergencystretch}{2em}
\multlinegap0pt
\usepackage{amssymb}
\usepackage{enumerate}
\usepackage[all]{xy}
\usepackage{tikz}
\newtheorem{lem}{Lemma}
 \newcommand{\bbox}{\mathbin{\Box}}
 \newcommand{\cc}{\mathbin{\textup{\copyright}}}
\title{Finite models for positive combinatorial and exponential algebra}
\author{Tumadhir Alsulami, Marcel Jackson}
\date{Source: arXiv:2411.05101v2 (CC BY 4.0)}
\begin{document}
\maketitle

\noindent\emph{Source and licence.} Finite models for positive combinatorial and exponential algebra. Version arXiv:2411.05101v2 (CC BY 4.0), distributed by arXiv under the Creative Commons Attribution 4.0 International licence, http://creativecommons.org/licenses/by/4.0/, as recorded in the arXiv licence metadata for this exact version and shown on the abstract page https://arxiv.org/abs/2411.05101v2. Full licence notice: \texttt{licenses/arxiv-2411.05101-CC-BY-4.0.txt}.\par\smallskip

\noindent\textit{Editorial scope.} Counted unit: the article's lem lem:gurevic (source0.tex lines 253--264). The article's complete original proof of it is delivered with it (source0.tex lines 265--296, one \texttt{proof} environment of 32 lines). No mathematics is written, repaired or re-derived: the statement and the proof are the article's own lines, in the article's own notation, and no retained line is substituted or re-flowed.  No further source-local context is needed: every cross-reference of every retained line resolves inside the two delivered spans.  Nothing was omitted from any retained span, so the package carries no omission record.  No retained line contains a citation, so the wrapper carries no bibliography and no bibliography entry.  No retained line contains a figure environment, an image-inclusion command or drawing code, so no image is delivered and the package declares no asset dependency.  Editorial formatting only, in addition: an AMS \texttt{amsart} preamble that restates the article's own declarations and macros used by the retained lines, namely the environments \texttt{lem} on separate counters, together with the standard packages those lines need.  The retained environments print in this document as Lemma 1 in order of appearance, whereas the article prints the same items with the numbering of its own full document.  The complete native source of the article as distributed in the arXiv e-print for this exact version is bundled unchanged as \texttt{sources/168184/source0.tex}.
\begin{lem}\label{lem:gurevic}
Let $\mathscr{L}= \{+,\cdot,\cc,!,\exp_2,0,1\}$ and $\Omega$ be the set of equations 
\begin{align*}
&\{1\cdot x\approx x\cdot 1\approx x,\  x\cc 0\approx 1\approx 0\cc x,\\
& 0+x\approx x+0\approx x,\ 0\cdot x\approx x\cdot 0\approx 0,\\
&0!\approx 1,\ 1!\approx 1,\ \exp_2(0)\approx 1,\ \exp_2(1)\approx 1+1\approx (1+1)!\}.
\end{align*}
Then there is a recursive function $B:\mathscr{L} \times \mathscr{L} \rightarrow \mathbb{N} $ such that, for any $\mathscr{L}$-terms $t_1, t_2$, and any set 
$\Sigma$ of valid equalities, the property
$\Sigma\cup\Omega\nvdash t_1\approx t_2$ holds
 if and only if $t_1\approx t_2$ fails in a model of $\Sigma\cup\Omega$, with cardinality at most $B(t_1,t_2)$.
\end{lem}

\begin{proof}
Throughout, $m$ will denote the number of variables in $t_1\approx t_2$.
Let $E= \Sigma\cup\Omega$ and $T_m$ be the set of $\mathscr{L}$-terms in variables $v_1,v_2, \dots ,v_m$ modulo equivalence of terms $s_1,s_2$ whenever $E \vdash s_1\approx s_2$.  In other words, $T_m$ is (isomorphic to) the $m$-generated free algebra in the variety defined by $E$, and so $t_1 \neq t_2$ in $T_m$ (or at least, the equivalence classes containing $t_1$ and $t_2$) means the same as $E\nvdash t_1\approx t_2$. 
Define a \emph{weight} function $w: \mathscr{L} \rightarrow N$ on terms by $w(0):=0 ,w(1):=1$, $w(v_i):=3$ (for each $i=1,\dots,m$) and by $w(t_1 \bbox t_2):=w(s_1)\bbox w(s_2)$ for $\bbox \in \{+,\cdot,\cc\}$, and $w(\bbox s):=\bbox(w(s))$ for $\bbox\in \{!,\exp_2\}$.  As the equalities in $E$ are valid in~$\mathbb{N}$ by assumption, it follows that $w$ is a well defined function $w: T_m \rightarrow N$.
  
Now inductively define a sequence $b_1(m),b_2(m),\dots$ as follows. Let $b_1(m)=2$, $b_2(m)=3$, $b_3(m)=5+m$ and let  $b_{k+1}(m)=3\cdot b_k(m)+3 \cdot (b_k(m))^2$ for $k\geq 3$.  
The intention is that $b_k(m)$ is an upper bound on the number of terms modulo $\Omega$ that have weight at most $k$.  
We now explain why this is so.   
For $b_1(m)=2$, observe that the laws in $\Omega$ and the definition of $w$ implies that the terms of weight at most $3$ are:
\[
\stackrel{\text{weight 1}}{\overbrace{0, 1}},\quad  \stackrel{\text{weight 2}}{\overbrace{1+1}},\quad  \stackrel{\text{weight 3}}{\overbrace{1+(1+1),(1+1)+1,v_1,\dots,v_m}}.
\]
For the inductive case, where $k\geq 3$, a term of weight at most $k+1$ either has weight at most $k$ (there are at most $b_k(m)$ of them), or has weight exactly $k+1$ and is of the form $s_1 \bbox s_2$ for some terms $s_1$, $s_2$ of weight at most $k$ (where $\bbox \in \{+,\cdot,\cc\}$),  or of the form $\bbox(s)$ for some term $s$ of weight at most $k$ (and where $\bbox\in \{!,\exp_2\}$).  
There are at most $3 \cdot (b_k(m))^2+2\cdot b_k(m)$ such possibilities, giving at most $3\cdot b_k(m)+3 \cdot (b_k(m))^2$ terms of weight at most $k+1$, as required.
Thus $\{s \in T_m \mid w(s)\leq j\}$ has at most $b_j(m)$ members. 
Let $K:=\max\{w(t_1),w(t_2),3\}$; we show that $B(t_1,t_2):=b_K(m)+1$ suffices.

 
Define an equivalence relation $\equiv$ on $T_m$ by $s_1\equiv s_2$ if $s_1$ and $s_2$ are already equivalent in $T_m$, or if $w(s_1), w(s_2)\geq K+1$.  
There are at most $B(t_1,t_2)=b_K(m)+1$ equivalence classes of $\equiv$ on $T_m$.   
The relation $\equiv$ is a congruence, as we now explain.  Assume $s\equiv t$.  If $s$ and $t$ are equivalent in $T_m$ then the stability of $\equiv$ holds trivially, so we assume that $w(s), w(t)\geq K+1$.  
For a unary operation $\bbox\in \{!,\exp_2\}$, the weight of $\bbox(s)$ and $\bbox(t)$ is strictly greater than $K$, because $\bbox$ is strictly increasing on inputs greater than $2$, and the weight of $s$ and $t$ is at least $3$.  
There are more cases when $\bbox$ is a binary operation in $\{+,\cdot,\cc\}$, but the idea is very similar.  
If~$u$ is a non-constant element of $T_m$, or a constant greater than $1$, then the weight of  $u\bbox s$, $s\bbox u$, $u\bbox t$, $t\bbox u$ is again greater than $K$, so that $u\bbox s\equiv u\bbox t$ and $s\bbox u\equiv t\bbox u$. 
When $u=1$, the weight also increases, except in the case of $\bbox=\cdot$.  
But $1\cdot s=s\cdot 1=s\cdot 1\equiv t=t\cdot 1$ anyway (this step uses $\Omega$), so again $\equiv$ is stable under translation by $u$.  
Stability similarly holds when $u=0$ and $\bbox$ is $+$.  
When $\bbox\in \{\cc,\cdot\}$, the laws in $\Omega$ ensure that $0\bbox s=s\bbox 0=c=0\bbox t=t\bbox 0$ for $c\in \{0,1\}$.  
So $\equiv$ is a congruence.  

To complete the proof, recall that $t_1\not\equiv t_2$, so lie distinct congruence classes in the quotient $T_m/{\equiv}$, which is then a  model of $E$ on at most $B(t_1,t_2)$-elements that fails $t_1\approx t_2$.  
\end{proof}

\end{document}
